\documentclass[10pt, a4paper]{article}
%-----------------------
%- 	PACKAGES & SETTINGS
%-----------------------
\usepackage[a4paper,top=3cm,bottom=4cm,left=4cm,right=4cm]{geometry}
\usepackage[utf8]{inputenc}
\usepackage[italian]{babel}
\usepackage{xcolor}
\usepackage{hyperref}
\hypersetup{
    colorlinks=true,
    filecolor=magenta,      
    urlcolor=darkgray,
    linkcolor=black
}
\urlstyle{same}
\usepackage{amsmath}
\usepackage{graphicx}
\graphicspath{ {images/} }
 
%-----------------------
%- 	TITLE
%-----------------------
\title{\textbf{Relazione Progetto OCaml PR 2}}
\author{\textbf{Venturi} Ludovico\\Docente: \href{http://pages.di.unipi.it/levi/}{Francesca Levi}}
\date{UNIPI, Gennaio 2020}


%-----------------------
%- 	DOCUMENT
%-----------------------
\begin{document}
%- 	INTRO
\pagenumbering{roman} 
\maketitle
\tableofcontents
\vfill
\begin{figure}[h]
	\centering
	\includegraphics[scale=0.3]{ocaml_logo}
	\label{fig:0}
\end{figure}

\clearpage

%- 	START DOC
\pagenumbering{arabic} 
\section{Scelte progettuali}
\begin{itemize}
\item tutti i valori associati alle chiavi in un dizionario hanno lo stesso tipo
\item i tipi assumibili dai valori associati alle chiavi sono solamente interi e booleani, nello specifico: \textit{Int(..), Bool(..)}
\item non sono ammessi valori \textit{Unbound} (corrispettivo di \textit{null})
\item per semplicità di progettazione si assume che il primo valore nel dizionario, al momento della creazione, definisca il tipo dei valori del dizionario
\item sono state introdotte le astrazioni funzionali multi argomento (funzioni con lista di parametri di dimensione variabile) per generalizzare dato che la \texttt{Fold} necessita di 2 argomenti
\end{itemize}
\subsection{Sintassi Astratta}
Estensione della sintassi astratta del linguaggio  \texttt{type exp = }\\
\texttt{...| CreateDict of (ide * exp)list | Insert of ide * exp * exp | Delete of  ide * exp | HasKey of ide * exp | Iterate of exp * exp	|
Fold of exp * exp | Filter of ide list * exp 	|.D..\\
...|FunCall of exp * exp list | FunArg of ide list * exp}\\\\
Estensione dei tipi esprimibili	\texttt{type evT = }\\
\texttt{...| DictVal of (ide * evT)list |...\\...| FunArgVal of ide list * exp * evT env }
\subsection{Runtime Type Checker}
Il \textit{Runtime Type Checker} è stato implementato per la creazione del dizionario, dove si verifica che tutti i valori siano consistenti nel tipo (o tutti interi, o tutti booleani). Influisce anche sull'operazione di inserimento, mentre sulle operazioni \texttt{Delete, HasKey, Filter} non ha alcuna influenza.\\ 
Per quanto riguarda le applicazioni di funzioni sui valori, ovvero \texttt{Fold, Iterate}, si demanda il type checking (già implementato nel linguaggio didattico) alle funzioni chiamate (sfruttando la \textbf{lazy evaluation} dei parametri): se viene chiamata la \texttt{Fold} di una funzione che opera su booleani su di un dizionario contenente valori interi al momento della prima chiamata di funzione viene generato un \textbf{type error}.
\section{Utilizzo}
Aprire l'interprete top-level di OCaml da terminale digitando \texttt{ocaml}.\\ Importare l'interprete del linguaggio qui discusso:
\begin{center}
\texttt{\emph{\#}\# use "venturi.ml";;}
\end{center}
e da qui valutare le espressioni, entrando nel vivo del \textbf{REPL}(ReadEvalPrintLoop).\\Vari testcase sono riportati nel file \texttt{testcase.ml}. \clearpage
\subsection{Esempi}
In \texttt{venturi.ml} vengono creati 4 ambienti:\\ \\
\texttt{
emptyenv Unbound;;(*in \emph{env0}*)\\{
.. "intDict" -> [("Birman", Eint(3));("Mainecoon", Eint(13));("Siamese", Eint(17));("Foldex", Eint(21))](*in \emph{env1}*)\\
.. "boolDict" -> [("Birman", Ebool(true));("Mainecoon", Ebool(false));("Siamese", Ebool(false));("Foldex", Ebool(false))](*in \emph{env2}*)\\
.. "emptyDict" -> [] (*in \emph{env3}*)}
}\\ \\
Sempre nello stesso file vengono riportate le prove \textbf{per ogni} operazione del dizionario, incluse tutte le eccezioni da esse generabili.\\
Alcuni esempi: 
\begin{itemize}
\item \textbf{Insert} 
\begin{itemize}
\item \texttt{
eval (Insert("Ragdoll", Eint(111), Den "intDict")) env3;;}\\\textit{- : evT =
DictVal
 [("Birman", Int 3); ("Mainecoon", Int 13); ("Siamese", Int 17);
  ("Foldex", Int 21); ("Ragdoll", Int 111)]}
\item \texttt{
		eval (Insert("Siamese", Eint(111), Den "intDict")) env3;;}\\\textit{Exception: InvalidArgumentException "key already present".
}
\end{itemize}
\item \textbf{Fold}
\begin{itemize}
\item \texttt{•} \\ \textit{•}

\item \texttt{•}\\ \textit{•}
\end{itemize}
\end{itemize}
\end{document}









